\section{Origin: how $2/9$ became visible in the first place}
Before any calculation, the question of where the central number of this
document comes from should be settled -- because it was not sought but found.
The route was the Hilbert-space translation of FFGFT. The bridge of
Docs.~230/231/232 maps the loop observables of the corpus to self-adjoint
operators on a fibre-valued Hilbert space; only when this translation was
carried through to the masses (Doc.~282) did the decisive structure appear:
the three lepton masses are the spectrum of a hermitian Z$_3$ circulant on
the $\mathbb{C}^3$ fibre of the full space
$\mathcal{H}=L^2(T^4)\otimes\mathbb{C}^3$. In the eigenvalues of this
operator the phase angle $\theta=2/9$ appears -- not as an ansatz one
inserts, but as a value the diagonalisation returns. $2/9$ is therefore
\emph{not} a starting point of the mass calculation but an order that became
visible only afterwards: it was the change of representation -- from the
$\xi$ ladder to the operator/spectral view -- that exposed it. This
counter-calculation therefore tests not an input but a \emph{discovery};
that is why the direction from the data (below) is meaningful at all. The
relation to the $\xi$ ladder, which computes the masses without $2/9$, is
clarified in the closing section \enquote{Two layers}.

\section{Purpose and direction of the calculation}
The FFGFT corpus fundamentally computes \emph{away from} $\xi=4/30000$: the
geometry yields relations that are subsequently checked against data. This
document reverses the direction. The starting point is exclusively the
empirical lepton masses (PDG), and the question is what the FFGFT relations
of the lepton sector make \emph{of the data} -- which structural numbers the
data return, with which uncertainties, and where residual tensions sit. The
mode is the compatibility mode: no new derivation, but an independent
counter-calculation per P35, with the level separation per P40 (only ratios
are exact; absolute values carry the bridge constants). All numbers in this
document come from \texttt{python/Dok292\_Skripte/leptonen\_empirie\_check.py}
(numpy-only, seed 20780458) and are reproducible with a single call.

The inputs are
\begin{align}
  m_e    &= 0.51099895069(16)\ \text{MeV},\\
  m_\mu  &= 105.6583755(23)\ \text{MeV},\\
  m_\tau &= 1776.86(12)\ \text{MeV}.
\end{align}
The $\tau$ mass is by orders of magnitude the least precise of the three; as
will become clear, it dominates all residual tensions.

\section{Level 1: the circulant -- Koide relation and angle}

\subsection{Koide ratio from the data}
The Koide ratio
$Q=(m_e+m_\mu+m_\tau)/(\sqrt{m_e}+\sqrt{m_\mu}+\sqrt{m_\tau})^2$ evaluated on
the PDG values gives
\begin{equation}
  Q = 0.666660511 \pm 0.0000068,
\end{equation}
where the uncertainty is carried almost entirely by $m_\tau$. The deviation
from the exact value $2/3$ is $-6.2\cdot10^{-6}$, i.e.\ $-0.91\,\sigma$. The
data are compatible with exactly $2/3$; they do not \emph{enforce} the exact
value, but they do not contradict it at $5$--$6$ significant digits.

\subsection{Circulant fit: amplitude and angle from three masses}
Following Doc.~282, the $\sqrt{m_k}$ are eigenvalues of a hermitian
Z$_3$ circulant, $\sqrt{m_k}=M\,(1+r\cos(\theta+2\pi k/3))$, with the two
pure structural numbers $r=\sqrt2$ and $\theta=2/9$ and the identity
$Q=(1+r^2/2)/3$. The Fourier extraction from the three data points yields
(assignment $k=0,1,2\to\tau,e,\mu$; all cyclic assignments give the same
angle)
\begin{equation}
  M = 313.84\ \text{MeV},\qquad
  r_{\text{emp}} = 1.4142005,\qquad
  \theta_{\text{emp}} = 0.22222963 \pm 8.4\cdot10^{-6}.
\end{equation}
The amplitude deviates from $\sqrt2=1.4142136$ by $1.3\cdot10^{-5}$ -- via
$Q=(1+r^2/2)/3$ this is the \emph{same} statement as the $Q$ deviation, not
an independent one. The angle lies $+0.89\,\sigma$ above
$2/9=0.22222222$, likewise $\tau$-limited.

\subsection{The $\tau$-free test: angle from $m_\mu/m_e$ alone}
Doc.~282 names $m_\mu/m_e$ as the sharpest test of whether $\theta$ is
exactly $2/9$: holding $r=\sqrt2$ fixed, $M$ cancels in the ratio, and the
two most precise masses in physics determine $\theta$ alone. Solving
$\bigl(1+\sqrt2\cos(\theta+4\pi/3)\bigr)^2/\bigl(1+\sqrt2\cos(\theta+2\pi/3)\bigr)^2
= m_\mu/m_e$ gives
\begin{equation}
  \theta(\mu/e) = 0.2222220471,\qquad |\theta-2/9| = 1.75\cdot10^{-7}.
\end{equation}
This independently reproduces the figure from Doc.~282
($\approx1.8\cdot10^{-7}$, seven significant digits). Two readings must be
kept apart. Measured against the tiny PDG uncertainty of $m_\mu/m_e$, the
residual is formally highly significant -- \emph{under the additional
assumption} that $r$ is exactly $\sqrt2$. If $r$ is allowed to float via
$Q$, the residual is not an independent tension but once again the same
single $\tau$ tension of $\sim0.9\,\sigma$. In the lepton sector there are
thus not four small deviations but \emph{one}, appearing in four guises.

\subsection{Forward check and the one prediction}
Conversely, setting $r=\sqrt2$ and $\theta=2/9$ exactly, the circulant
reproduces $m_\mu/m_e$ to $+0.0010\,\%$ and $m_\tau/m_\mu$ to
$+0.0060\,\%$. Inverting for $m_\tau$ turns this into a falsifiable forward
statement: from $(m_e, m_\mu)$ and the two structural numbers it follows that
\begin{equation}
  m_\tau^{\text{pred}} = 1776.968\ \text{MeV}
\end{equation}
(the three inversion routes -- $Q=2/3$ exact, $\theta=2/9$ exact in the
three-mass fit, direct prediction from $(m_e,m_\mu;\sqrt2,2/9)$ -- give
$1776.969$, $1776.967$ and $1776.968$ MeV, i.e.\ the same value). Against
PDG $1776.86\pm0.12$ MeV this is $+0.90\,\sigma$. A $\tau$-mass measurement
at the targeted Belle-II level of $\sim0.05$ MeV decides this tension: if
the value lands near $\approx1776.97$, the $\tau$ direction is resolved and
the pair $(\sqrt2, 2/9)$ is confirmed at the $10^{-5}$ level; if it lands
significantly below, the approximation at this level is refuted. This is a
genuine prediction in the sense of the corpus methodology, not a
retrodiction: the value is fixed before the measurement and is booked here
in advance. What the test does \emph{not} decide is shown by the precision
analysis in Section~\ref{sec:precision}.

\section{Level 2: the $\xi$ ladder -- ratios and reverse $\xi$}

\subsection{Ladder ratios against the data}
The ladder representation $m_i = r_i\,\xi^{p_i}\,v$ (Docs.~006/046; leptons:
$r_e=4/3$, $p_e=3/2$; $r_\mu=16/5$, $p_\mu=1$; $r_\tau=25/9$, $p_\tau=2/3$
per K2) makes the ratios parameter-free:
\begin{align}
  \frac{m_\mu}{m_e}\Big|_{\text{ladder}} &= \tfrac{12}{5}\,\xi^{-1/2} = 207.85
  \quad(\text{PDG } 206.77;\ +0.52\,\%),\\
  \frac{m_\tau}{m_\mu}\Big|_{\text{ladder}} &= \tfrac{125}{144}\,\xi^{-1/3} = 16.992
  \quad(\text{PDG } 16.817;\ +1.04\,\%),\\
  \frac{m_\tau}{m_e}\Big|_{\text{ladder}} &= \tfrac{25}{12}\,\xi^{-5/6} = 3532
  \quad(\text{PDG } 3477;\ +1.57\,\%).
\end{align}
The ladder thus carries the orders of magnitude and the $\sim1\,\%$
accuracy -- not the precision of the circulant. Both levels are to be read
separately, as declared in Doc.~282 and P40: the circulant delivers four to
five digits from $(\sqrt2, 2/9)$, the ladder delivers the hierarchy from
$\xi$.

\subsection{Reverse $\xi$: what the data return for $\xi$}
Solving the three ladder ratios for $\xi$, the data return
\begin{equation}
  \xi_{\mu/e} = 1.347\cdot10^{-4},\qquad
  \xi_{\tau/\mu} = 1.375\cdot10^{-4},\qquad
  \xi_{\tau/e} = 1.358\cdot10^{-4},
\end{equation}
i.e.\ $+1.1$ to $+3.2\,\%$ above $4/30000=1.3333\cdot10^{-4}$. The three
values scatter among each other by $\sim2\,\%$ -- the ladder is consistent
in the ratios at the percent level, not exact. This is the honest empirical
status of the ladder representation, unchanged from the declaration in
Doc.~006 (mean deviation below $1\,\%$ at the mass level; correspondingly
$1$--$3\,\%$ at the ratio level).

\subsection{Bridge constants per P40}
Per lepton, solving for the bridge constant
$v_{\text{eff}}=m_i/(r_i\,\xi^{p_i})$ yields
\begin{equation}
  v_{\text{eff}}^{(e)} = 248.9\ \text{GeV},\qquad
  v_{\text{eff}}^{(\mu)} = 247.6\ \text{GeV},\qquad
  v_{\text{eff}}^{(\tau)} = 245.1\ \text{GeV},
\end{equation}
scattered around the Higgs vacuum expectation value $246.22$ GeV
($+1.10/+0.58/-0.46\,\%$). Equivalently: the required $r_{\text{eff}}$ at
fixed $v$ deviate by the same percentages from the declared rational values
($r_e$: $1.348$ instead of $4/3$, etc.). This is precisely the P40 finding
seen from the data side: the rational structural numbers of the ladder hold
at the percent level; the residual deviation is absorbed into the bridge
constant and is \emph{not} fully cancelled in ratios.

\subsection{$K_{\text{frak}}$ cross-check: no simple pattern}
The obvious question of whether the ladder residuals follow the fractal
correction factor is answered in the negative by the calculation:
$K_{\text{frak}}=1-100\xi$ corresponds to $-1.333\,\%$, the half and
three-halves powers to $-0.665$ and $-1.995\,\%$ -- the measured deviations
$+0.52$ and $+1.04\,\%$ fit no single $K_{\text{frak}}$ power in either
magnitude or sign. The only striking feature is that the $\tau/\mu$
deviation is quite precisely twice the $\mu/e$ deviation; whether structure
or coincidence lies behind this remains unbooked per P35. The point is
declared open, not explained away.

\subsection{Cross-ladder coupling: a generation-linear correction law}
Before this finding is booked, the booking level must be fixed, since it
decides admissibility. There are two \emph{mutually exclusive} calculation
modes. \emph{Mode~1} computes all lepton ratios purely from $\xi$, without a
reference point: then $\mu/e$, $\tau/\mu$ and $\tau/e$ are all genuine
predictions (each $0.5$--$1.6\,\%$ off the data), and all three are
legitimate test points. \emph{Mode~2} applies P42: $m_\mu/m_e$ is then the
declared anchor, \emph{consumed} and no longer a test quantity -- one may no
longer book its ladder deviation, else the same quantity is used twice (once
as anchor, once as test). The coupling analysis below belongs
\emph{exclusively} to Mode~1, the reference-free one; in Mode~2 it is void,
because there $\mu/e$ drops out of the comparison and the only independent
ladder test left is $m_\tau$ -- which in any case runs through the circulant,
not the coarse ladder. Within Mode~1: the statement \enquote{no simple
pattern} concerns a \emph{single} $K_{\text{frak}}$ power. Asking instead whether the ladder stages are
coupled \emph{to each other} does reveal structure. Writing the required
correction per place as an effective factor $N_g$ via $K=1-N_g\,\xi$ (or as
an exponent $p$ via $K_{\text{frak}}^{p}$), the data give
\begin{align}
  \text{$\mu/e$:}\quad & N_1 = 38.89,\quad p_1 = 0.3873,\\
  \text{$\tau/\mu$:}\quad & N_2 = 77.06,\quad p_2 = 0.7694,\\
  \text{$\tau/e$:}\quad & N_3 = 115.55,\quad p_3 = 1.1567.
\end{align}
Two relations here are not accidental. First, the exponent structure is
\emph{multiplicatively consistent}: $p_3 = p_1+p_2$ to $6\cdot10^{-14}$. The
$\tau/e$ place therefore carries \emph{no} new information -- it is the
product of the two stages, as a throughgoing multiplicative correction
$K^{p_1}\cdot K^{p_2}=K^{p_1+p_2}$ requires. Second, the effective factors
stand in the ratio $N_1:N_2:N_3 = 1:1.981:2.971 \approx 1:2:3$: the
correction grows \emph{linearly with generation number},
$N_g \approx g\cdot N_0$ with $N_0 \approx 38.6$ (the three individual
estimates $N_1$, $N_2/2$, $N_3/3$ are $38.89$, $38.53$, $38.52$ --
consistent among themselves to below $1\,\%$).

This is the structural reason why a constant $K_{\text{frak}}$ does not
close the ladder, and why asking for the \enquote{right} number instead of
$100$ (such as $89$) misses the point: the ladder demands not a fixed factor
but a \emph{generation-counting} one, $N_g=g\cdot N_0$. At $\alpha$ only one
scale is in play (the $\mu e$ mean) -- one factor suffices, and $\alpha$
fits. In the ladder three generations superpose, so $N$ grows linearly; a
constant factor cannot do this in principle. The generation-dependent
dressing previously named as an obligation thereby acquires a concrete
form: to first approximation it is linear in $g$.

The tolerance balance shows how far the law reaches (reference-free Mode~1,
common $N_0=38.65$):
\begin{center}
\begin{tabular}{lrrrr}
\hline
Place & raw & raw [$\sigma_{\text{exp}}$] & corrected & corr.\ [$\sigma_{\text{exp}}$]\\
\hline
$\mu/e$    & $+0.521\,\%$ & $2.4\cdot10^{5}$ & $+0.0033\,\%$ & $1.5\cdot10^{3}$\\
$\tau/\mu$ & $+1.038\,\%$ & $154$            & $-0.0031\,\%$ & $0.47$\\
$\tau/e$   & $+1.565\,\%$ & $232$            & $-0.0052\,\%$ & $0.78$\\
\hline
\end{tabular}
\end{center}
The law reduces the ladder deviation by about a factor of $300$, from
$0.5$--$1.6\,\%$ to $\sim5\cdot10^{-3}\,\%$. The two $\tau$ places thereby
fall \emph{within} their measurement tolerance ($<1\,\sigma$) -- for them
the law is practically complete. The $\mu/e$ entry stays at
$\sim1.5\cdot10^{3}\,\sigma$, not because the correction fails there but
because $\mu/e$ is measured about $500$ times more precisely: even a
$0.003\,\%$ residual is enormous in $\sigma$ units. This is exactly where
the mode boundary bites: this $\mu/e$ entry is a statement \emph{only} in
reference-free Mode~1; under P42, $\mu/e$ is the anchor and drops out of the
balance. The remaining residual of the two $\tau$ places sits in the fine
deviation of the ratio from exact $1:2:3$ ($\sim1$--$2\,\%$) and in the
still-undetermined base unit $N_0$ -- that is the next open calculation, not
a contradiction.

Open -- and per P35 expressly not explained away -- are two things. The base
unit $N_0\approx38.6$ is not yet derived from the geometry; obvious
candidates ($100/e=36.8$, $100/\varphi^2=38.2$, plain $39$) do not match it
closely enough to be claimed. And the phase sector -- the $\delta$ fine
value from Doc.~291 with $p_\delta=-0.88$ -- does \emph{not} obviously fit
the same generation-linear magnitude ladder; whether the magnitude and
phase sectors couple through a common law is a separate question not
answered here. The finding is thus not the solution of the ladder residual
but its sharpening: \enquote{three crooked exponents} becomes \enquote{one
generation-linear law $N_g=g\cdot N_0$ with $N_0$ still undetermined}.

The limit of this claim, following from the mode separation, must be stated
explicitly. The $1:2:3$ pattern is supported at three points \emph{only} in
the reference-free Mode~1. As soon as P42 is applied and $m_\mu/m_e$ is
declared the anchor (the level needed for the circulant precision and the
$m_\tau$ prediction), $\mu/e$ drops out as a test point, and the law rests
only on $\tau/\mu$ and $\tau/e$ -- which share the same $\tau$ mass and whose
precision is low. The apparently threefold confirmation is thus tied to
reference-freeness. Both at once -- $\mu/e$ as anchor \emph{and} as an
$N_g$ test point -- is inadmissible double-booking. The circulant, which
already carries $\mu/e$ to $10^{-5}$ (route~A in
Section~\ref{sec:precision}), does not need the $g$ correction anyway; it is
a statement about the coarse $\xi$ ladder in the reference-free mode alone,
not about the precise circulant layer.

A common misunderstanding must be cleared up here: the $100$ in
$K_{\text{frak}}=1-100\,\xi$ and the base unit $N_0\approx38.6$ are
\emph{not} the same quantity, so the ladder by no means needs \enquote{fewer
recursions} than the $\alpha$ sector. The $100$ is the single fixed scale of
a corpus-wide constant (X-flip, $T_{\text{CMB}}$, $E_0$ bridge); $N_0$ is the
factor \emph{per generation} in a different sector. The difference is not
\emph{fewer} but \emph{other}: not a fixed factor but a generation-counting
$N_g=g\cdot N_0$ -- already $77$ at $g=2$, $116$ at $g=3$, i.e.\ partly
\emph{more} than $100$. Comparing the two numbers as \enquote{recursion
depths} therefore misses the point.

In assessing possible expressions for $N_0$ the correct measure is not
point-exact agreement -- that is unreachable in principle -- but consistency
\emph{within the tolerance} to which $N_0$ can be determined at all. That
tolerance is considerable: in the referenced mode $N_0$ follows from the
$\tau$ places alone, and their precision is limited by the $\tau$ mass
($\pm0.12$ MeV). Propagating this gives $N_0 = 38.52 \pm 0.21$, a relative
determination uncertainty of about $0.5\,\%$. A candidate expression is
therefore to be judged not by whether it hits $38.6$ to the decimal, but by
whether it falls into this interval.

By that measure -- and still as an \emph{unproven} observation per P35, not
a derivation -- $100/\varphi^2 = 38.20$ is an \emph{admissible} candidate:
its distance to $N_0$ is $1.6\,\sigma_{N_0}$, i.e.\ within the tolerance. As
a point value the expression misses $38.6$; within the tolerance to which
$N_0$ is known, it matches. Were $N_0=100/\varphi^2$, the ladder base unit
would follow from the $K_{\text{frak}}$ scale through a single factor
$\varphi^{-2}$ -- magnitude and correction sectors hanging on the same
constant. The finding is thus not \enquote{narrowly missed and discarded}
but \enquote{compatible with the present determination, not refuted}. A
sharper $\tau$ mass (Belle~II) or a geometric derivation of $N_0$ will decide
between $100/\varphi^2$ and the alternatives; until then it remains an
admissible, not a proven, expression.

\section{Relation to the measurement precisions}
\label{sec:precision}
The $\sigma$ statements so far refer to one quantity each. The full
question is: how do the FFGFT structural numbers sit relative to the
\emph{entire} precision structure of the data? A two-parameter analysis
answers it. The two independent ratios $R_1=m_\mu/m_e$ and
$R_2=m_\tau/m_\mu$ determine the pair $(r,\theta)$ exactly; their relative
measurement precisions, however, differ dramatically:
$\sigma_{R_1}/R_1 = 2.2\cdot10^{-8}$ versus
$\sigma_{R_2}/R_2 = 6.8\cdot10^{-5}$ -- the $\mu/e$ ratio is measured
about $3100$ times more precisely. Error propagation onto $(r,\theta)$
therefore yields an extremely anisotropic covariance: a \emph{thin}
principal direction with $\sigma \approx 3.4\cdot10^{-10}$ (carried by
$\mu/e$) and a \emph{thick} one with $\sigma \approx 1.7\cdot10^{-5}$
(carried by $\tau$); the correlation between $r$ and $\theta$ is
practically $-1$.

In this picture the point $(\sqrt2, 2/9)$ lies away from the data fit
$(1.41420051,\ 0.22222963)$ by $0.90\,\sigma$ in the thick direction --
this is the single $\tau$ tension of Section~2 -- and by about
$450\,\sigma$ in the thin one. Spelled out: the circulant with
\emph{exactly} $(\sqrt2, 2/9)$ misses $m_\mu/m_e$ by $1.0\cdot10^{-5}$
relative, while the measurement stands at $2.2\cdot10^{-8}$. As an exact
pair the two structural numbers are thus already excluded by today's data;
they are confirmed at the $10^{-5}$ level -- precisely the \enquote{four to
five digits} declared in Doc.~282, here made precise as a measurement
relation. This also bounds the reach of the $m_\tau$ prediction: a perfect
Belle-II result resolves only the thick $\tau$ direction; the residual of
$1.75\cdot10^{-7}$ in $\theta$ (equivalently $10^{-5}$ in $R_1$) is an
\emph{already measured}, standing offset that any exactness claim would
have to explain -- radiative corrections to the pole masses are the obvious
candidate, though their typical size ($\alpha/\pi \sim 10^{-3}$) rather
deepens the puzzle of why the agreement reaches $10^{-5}$ at all. The point
is booked open per P35.

\subsection{Thought model: inverting the precision statement}
The $450\,\sigma$ finding rests entirely on trust in the quoted precision
of $m_\mu/m_e$. The inversion -- the precision might not mean what it
claims -- is admissible as a thought model and comes in three strengths.
\emph{Weak:} the experimental uncertainty would be underestimated by a
factor of $\sim450$ -- quantitatively hardly tenable, since the same
precision hangs in a web of independent cross-checks (muonium spectroscopy,
$g\!-\!2$, QED tests). \emph{Intermediate:} $m_\mu/m_e$ is not a
read-off but a quantity \emph{extracted} through bound-state QED; the
quoted precision presupposes the completeness of the extraction theory, and
a missing term at $10^{-5}$ would dissolve the conflict -- $450$ times the
estimated theory error, borderline, but logically open. \emph{Strong (and
FFGFT-native):} the precision is real but applies to the \emph{dressed}
quantity -- the pole masses in the SI frame -- while the circulant relation
would live on the geometric level beneath. Then nothing contradicts
anything: the $10^{-5}$ offset would not be an error but the
\emph{measured thickness of the dressing layer} in the $\mu/e$ channel.
This is the consistent extension of P40: in ratios only the
generation-\emph{independent} part of the fractal/SI correction cancels; a
generation-dependent remainder survives -- and would have exactly the shape
of this offset. The price of the strong reading is an obligation: the
offset would have to be computable from the dressing
(generation-dependent radiative structure), otherwise the reading is
unfalsifiable. It is booked here as a thought model, not as a resolution;
the contradiction finding of the previous section stands for the pole-mass
level.

\subsection{Resolved by declaration: the reference point $m_\mu/m_e$ (P42)}
The framework decides the question not through one of the thought models but
through the claim layer: FFGFT fundamentally keeps the purely theoretical
derivation from $\xi$, and the lepton sector receives -- analogous to the
cosmic scale (R39) -- exactly \emph{one} declared reference point, the
ratio $m_\mu/m_e$. At structural $r=\sqrt2$ (from $Q=2/3$) it fixes the
operative angle $\theta_{\text{ref}}=0.2222220471$; the structural number
$2/9$ is its rational address to seven digits, not an exactness claim at
the pole-mass level. The $450\,\sigma$ finding is thereby no longer a
falsification but part of the calibration. Two delimitations: the status of
$\xi$ is untouched (FFGFT claims $\xi$ as forced by the geometry; that
dispute is not adjudicated here), and the lepton masses are in themselves
not an input of the theory -- the reference point belongs to the comparison
level, it locates the geometric relation in the pole-mass data and would
not be needed without the intent to compare. The honest price lies in the
test balance: $m_\mu/m_e$ is consumed as anchor and drops out as a test
quantity; the test quantity is $m_\tau$ alone. The gain is an
input-error-free sharp prediction:
\begin{equation}
  m_\tau = 1776.9690 \pm 0.0001\ \text{MeV},
\end{equation}
where the uncertainty stems solely from $(m_e, m_\mu)$ and practically
vanishes; against PDG the single $+0.91\,\sigma$ tension remains, and the
Belle-II world average decides with no escape route. For the bridge
question (Doc.~291) the declaration sharpens the target: the dynamics must
select $\theta_{\text{ref}}$; $2/9$ remains its address and the
pre-registered Stage-10 target, whose orbit resolution ($\sim0.1$) is
unaffected by the $10^{-7}$ difference. Booked as P42 in Doc.~190.

The same view calibrates the $\xi$ ladder: its residuals correspond to
$2.4\cdot10^{5}$ ($\mu/e$), $154$ ($\tau/\mu$) and $232$ ($\tau/e$)
experimental standard deviations. The ladder residuals are therefore not
data-limited but theory-side (P40 bridge constants); measurement precision
has outrun the ladder representation by two to five orders of magnitude.
Finally, the test power of the $m_\tau$ prediction (separation of
prediction from today's central value) at
$\sigma_{m_\tau}=0.12/0.05/0.02$ MeV is $0.9/2.2/5.4\,\sigma$
respectively -- a single Belle-II measurement is indicative; only the
combined world average is decisive.

\section{The fine-structure constant from the lepton masses}
The fine-structure constant belongs in this counter-calculation because in
FFGFT it is a lepton-mass statement: Doc.~011 sets
$\alpha = \xi\,(E_0/\text{MeV})^2$ with the structural energy scale
$E_0=\sqrt{m_e\,m_\mu}$ and the declared fractal correction,
$\alpha^{-1} = (7500/E_0^2)\,K_{\text{frak}}$. From the data side: the
measured $\alpha$ requires $E_0^{\text{req}}=7.39798$ MeV, the geometric
mean gives $7.34788$ MeV -- ratio $1.006819$, against
$K_{\text{frak}}^{-1/2}=1.006734$. The declared correction factor thus
closes the gap up to $8\cdot10^{-5}$; the closed, entirely parameter-free
relation
\begin{equation}
  \alpha \;=\; \frac{\xi\, m_e\, m_\mu}{K_{\text{frak}}\ \text{MeV}^2},
  \qquad K_{\text{frak}}=1-100\,\xi,
\end{equation}
gives $1/\alpha = 137.059$ against measured $137.035999$.

The decisive point here is that $\alpha$ does \emph{not} run through a
fitted bridge constant but through \emph{two independent routes} to the
characteristic energy $E_0$ that confirm each other (Doc.~011). Route~1 is
empirical: $E_0=\sqrt{m_e m_\mu}=7.348$ MeV, the geometric mean of the
measured masses. Route~2 is independently geometric:
$E_0^2=4\sqrt2\,m_\mu/\xi^{p}$ constructs $E_0=7.398$ MeV from $\xi$ and
$m_\mu$ alone, without using the geometric mean -- not a back-calculation
from $\alpha$. Two consistencies that know nothing of each other coincide:
(i) route~2 hits $7.398$ and thereby yields
$\alpha=\xi\,(E_0/\text{MeV})^2$ to $+5\cdot10^{-6}$; (ii) the difference
of the two routes, $7.398/7.348=1.00682$, equals $K_{\text{frak}}^{-1/2}$
to $9\cdot10^{-5}$. The $\alpha$ point is therefore \emph{not} a
calibration but overdetermined: two separate derivations meet to
$8\cdot10^{-5}$, and the recursion depth of the correction is not free but
structurally fixed -- $E_0$ enters $\alpha$ quadratically, so the
$K_{\text{frak}}^{1/2}$ bridge at the $E_0$ level corresponds to a full
$K_{\text{frak}}$ order at the $\alpha$ level. That the \emph{same} factor
$K_{\text{frak}}=1-100\,\xi$ which in the corpus carries the damped X-flip
and the $T_{\text{CMB}}$ correction closes exactly the gap between the two
$E_0$ routes is a cross-sector consistency, not a fit.

\subsection{The $\alpha$ residual as a second witness to the dressing layer}
The overdetermination of $\alpha$ permits a statement that the $\theta$
case (Section~\ref{sec:precision}) only hints at. Taking both anchors as
exact -- the $\alpha$ measured to $1.5\cdot10^{-10}$ and the purely
geometric route-2 construction $E_0^2=4\sqrt2\,m_\mu/\xi^{p}$, which
carries no measurement error -- the only quantity with any freedom in
$\alpha=\xi\,E_0^2$ is the empirical route~1, i.e.\ $E_0=\sqrt{m_e m_\mu}$
and hence the pole masses themselves. The $K_{\text{frak}}$ bridge closes
$99.99\,\%$ of the gap between route~1 and the $E_0$ required by $\alpha$;
a residual of $8.4\cdot10^{-5}$ remains. Measured against the precision of
the masses ($1.1\cdot10^{-8}$ at the $E_0$ level) this residual is
$\sim7700\,\sigma$ large -- too large to be mass measurement noise, and
therefore a statement \emph{about} the pole masses: under the strong
reading (route~2 exact), $\sqrt{m_e m_\mu}$ is not the geometrically exact
quantity.

Crucially, this $\alpha$ residual holds \emph{mode-independently} -- and
precisely because of the two routes. Route~2 does \emph{not} use $m_e$: it
constructs $E_0$ from $\xi$ and $m_\mu$ alone. Hence $m_\mu/m_e$ is
\emph{not} an anchor here; $m_e$ is nowhere inserted and then recovered. The
agreement of route~2 (without $m_e$) with route~1 (with $m_e$) is therefore a
genuine prediction about $m_e$, and the residual of $8\cdot10^{-5}$ is not
self-reference but a genuine test of the pole mass against the geometry -- it
remains a statement whether one computes reference-free or with P42.

This allows a precise placement -- sharper than a symmetric \enquote{two
witnesses} phrasing would give. The comparison with the $\theta$ side
(Section~\ref{sec:precision}) reveals an \emph{asymmetry}: there
$\theta_{\text{ref}}$ is obtained from $m_\mu/m_e$, so $\mu/e$ is the
anchor, and the $\sim450\,\sigma$ in the $\mu/e$ channel are half
self-reference -- a witness \emph{only} in the reference-free mode. The
$\alpha$ residual, by contrast, is \emph{mode-independent}, because route~2
does not use the second anchor ($m_e$) at all. The strong evidence, valid in
every mode, for a dressing layer between the geometrically fundamental
quantity and the QED-dressed pole mass is therefore $\alpha$ alone; the
$\theta$ side supports the same conclusion, but only under
reference-freeness. The honest price remains: the statement presupposes the
strong reading (route~2 as exact), and the obligation -- to \emph{compute}
the dressing rather than name it as a residual -- is open. As a calibration
artefact, however, the $\alpha$ residual cannot be dismissed: $\alpha$ is too
overdetermined and the mass precision too high for that.

Conceptually -- almost without computational consequence, yet central to
the architecture -- the following holds: in the extended natural unit
system $\alpha=1$ (Docs.~011/014, charge dimensionless); the SI value
$1/137.036$ is bookkeeping of the unit translation, not a fundamental
coupling. The same architecture carries gravitation:
$G=\xi^2/(4 m_{\text{char}})$ times conversion (Docs.~012/013) makes $G$
a bookkeeping quantity relative to SI -- reproducing the SI value is
calibration per R40; only the forward chain is testable. The structural
gain lies elsewhere: a coupling that is only bookkeeping need not be
quantised -- the quantum-gravity problem, on which most extensions of the
Standard Model founder, does not arise in this architecture.

\section{Two layers that do not need each other: what $2/9$ provides and what the SI conversion provides}
Two questions that are easily conflated must be separated here.

\emph{First: is $2/9$ needed to compute the masses?} No. The historical and
still valid route is the $\xi$ ladder (Docs.~006/046): $m_i=r_i\,\xi^{p_i}\,v$
with $r\in\{4/3,16/5,25/9\}$ and $p\in\{3/2,1,2/3\}$ yields the masses from
$\xi$ and the ladder numbers alone, with $2/9$ entering nowhere. The
structural number $2/9$ appears only from Doc.~282 onward and in a different
role: not as an input generating the masses, but as an \emph{order
discovered within them} -- the three $\sqrt{m_k}$ are eigenvalues of a
Z$_3$ circulant with phase angle $2/9$. $2/9$ replaces the previously fitted
ladder coefficient by a pure number; it is compression and explanation, not
a prerequisite. One could delete $2/9$ from the mass sector and the ladder
would keep computing -- what would be lost is the \emph{insight} into why
$Q$ is exactly $2/3$, not the numbers. $2/9$ thus belongs to the
explanatory layer, not the computational one.

\emph{Second: does this counter-calculation confirm that the SI conversion
is necessary and correct?} In two separate degrees. As \emph{necessary} it
confirms clearly: every absolute number ($m_\tau$ in MeV, $\alpha=1/137$,
$G$) requires the bridge constant to reach SI, while the ratios
($m_\mu/m_e$) do not. This two-level structure is exactly P40, and the
$450\,\sigma$ finding is its sharpest evidence: it shows that the measured
pole-mass level (SI-defined, QED-dressed) is a \emph{different} level from
the geometric one -- without the conversion step the difference would not
exist. As \emph{correct in the sense of validated to measurement
precision}, however, it does \emph{not} confirm: the conversion closes
about $98.8\,\%$ of the gap at $\alpha$ (through the same factor
$K_{\text{frak}}$ that also operates in other sectors), but the remaining
$\sim10^{-4}$ stays open. The SI conversion is therefore structurally
necessary and consistent across sectors, but it is calibration, not a test
passed to the decimal.

The real finding lies in the independence: the ladder computes masses
without $2/9$, the circulant explains $Q$ without the ladder, the SI
conversion translates without $2/9$. That all three nonetheless share the
same $\xi$ and the same factor $K_{\text{frak}}$ is the consistency of the
framework -- not a dependence of one layer on another.

\section{Placement}
The counter-calculation from the data side confirms the two-level
architecture of the corpus without adding to it or sparing it anything. At
the circulant level the data return the two structural numbers at the
$10^{-5}$ level; the precision analysis separates two findings: in the
$\tau$-limited (thick) direction a single $0.9\,\sigma$ tension, carrying
the pre-booked prediction $m_\tau = 1776.968$ MeV -- and in the
$\mu/e$-precise (thin) direction an already measured $10^{-5}$ offset that
excludes the \emph{exactness} of the pair $(\sqrt2, 2/9)$ and is booked as
open fine structure. At the ladder level
the data return $\xi$ at $1$--$3\,\%$, with honestly open residuals showing
no recognisable $K_{\text{frak}}$ pattern. What this calculation does
\emph{not} show: no derivation of $\theta=2/9$ (its dynamical locus remains
the open question of the bridge work), no sharpening of the ladder, no new
physics. It shows that the declared claims of both levels are independently
reproducible starting from the PDG data -- in the direction in which nothing
can be fitted.

\subsection{The residual as the signature of incompleteness}
The counter-calculation has sharpened the statements -- from \enquote{two
pure numbers, no fit} to measured quantities with known decimals and known
limits. Yet at every level a residual remains: the $10^{-7}$ offset in
$\theta$, the $1$--$3\,\%$ of the ladder, the open $\delta$ fine value
(Doc.~291), the $K_{\text{frak}}$ pattern that does not close across the
ladder stages. (The $\alpha$ point does \emph{not} belong here: it is
overdetermined via two independent $E_0$ routes and its recursion depth is
structurally fixed -- see above.) None of these residuals is a \emph{contradiction} -- the theory
holds everywhere to its respective precision -- but none is \emph{zero}. A
residual that shrinks the closer one looks without vanishing is the
signature of an approximation to something deeper, not of a completed
foundation. The residual is therefore not an error but the honest evidence
that the theory is \emph{not complete}.

What is remarkable is \emph{where} it sits. The geometric core carries no
residual: $Q=2/3$ is purely geometric and practically exact; in the extended
natural system $\alpha=1$; the ladder ratios are rational numbers. The
residual arises throughout \emph{only} at the transition into the measured:
the $450\,\sigma$ break only at the pole-mass level, the ladder deviation
only in the bridge constant. The $\alpha$ conversion is the exception that
sharpens the rule: there \emph{two} independent $E_0$ routes fix the
recursion depth, so the translation is not free -- precisely why $\alpha$
is the one place without an open residual in this sense. Where only
\emph{one} route exists (the ladder ratios), the depth of the correction
stays undetermined, and that is where the residual sits.
The residual thus lives not in the geometric core but in the
\emph{translation layer} -- systematically, not randomly distributed. This
localises the incompleteness precisely: not in $\xi$, not in the circulant,
but in the not-yet-\emph{derived} fractal or QED dressing. The theory would
be complete only once it \emph{computes} this residual instead of absorbing
it into a bridge constant -- once $K_{\text{frak}}$ not merely fits but
follows generation-dependently from the geometry. The measure for this is
expressly \emph{not} a point-exact landing on the last decimal: a geometric
derivation is not meant to hit the QED-dressed pole mass digit by digit --
that would be neither possible nor the goal -- but to fix the underlying
structure far enough that the predictions stay \emph{within the tolerances}
to which the quantities can be determined at all. Completeness here means: no
residual left that protrudes \emph{beyond} those tolerances and is closed
only by a free parameter. Until then, each such residual is the pointer to
the next open derivation, against which completeness -- in the sense of
tolerance, not of the last digit -- will be measured.
